% Deutsche Parallelfassung zu foundations/information_modelling.tex

\chapter{Informationsmodellierung}
\label{chap:information-modelling}

\section{Ein Alignment, mehrere Darstellungen}

Das bearbeitete Alignment erscheint nun an mehreren Stellen. Eine GIS-Sicht
verortet es zwischen Gelände und Netzen. Ein CAD-Modell legt die genaue
Konstruktionsgeometrie offen. Ein BIM- oder Austauschmodell verbindet es mit
Projektobjekten und Lieferabläufen. Eine Vermessung hält Beobachtungen des
physischen Gleises fest. Eine Berechnung bewertet eine Folge, und ein
Entscheidungsnachweis hält fest, was ein verantwortlicher Ingenieur annimmt.

Das bloße Zusammenführen dieser Aufzeichnungen ergibt noch kein kohärentes
Informationsmodell. Dieselbe Koordinate kann ein Entwurfsziel, einen physisch
realisierten Zustand oder eine gemessene Schätzung ausdrücken. Derselbe
Identifier kann zwischen Dateien kopiert werden, ohne zu beweisen, dass sich
die Aufzeichnungen auf dasselbe Engineering Object beziehen. Eine Regel kann
neben einem Modell gespeichert sein, ohne für die aktuelle Aufgabe anwendbar
zu sein.

Informationsmodellierung beginnt hier: nicht mit einem größeren Container,
sondern mit einer Darstellung dessen, worüber die Information eine Aussage
macht, welche Rolle jede Aufzeichnung spielt und welche Beziehungen ihre
Nutzung rechtfertigen. Auch die Interoperabilitätsforschung unterscheidet den
Datenaustausch von den semantischen, organisatorischen und menschlichen
Bedingungen, die für Zusammenarbeit erforderlich sind
\cite{Turk2020Interoperability}.

\section{Was Repräsentationen beitragen}

Repräsentationen sind unverzichtbar, weil Ingenieurarbeit Auswahl erfordert.
Geodatenstandards ordnen Zugriff, Verarbeitung, Visualisierung und
Erschließung räumlicher Informationen \cite{OGCStandardsProgram2025}.
CAD-Geometrie unterstützt genaue Konstruktion und technische Kommunikation
\cite{KimKumarSnider2015Geometry}. IFC stellt standardisierte
Beschreibungen für den Austausch von Informationen über gebaute Anlagen
zwischen Anwendungen und Beteiligten bereit; Modellsichten werden auf
anerkannte Arbeitsabläufe zugeschnitten
\cite{BuildingSMARTIFC4x3Scope}.

Jede Repräsentation beantwortet einen Zweck. Diese Stärke kennzeichnet zugleich
ihre Grenze. Keine Kodierung, Modellsicht, Zeichnung oder Datenbankzeile
begründet für sich genommen:

\begin{itemize}
\item ob zwei Aufzeichnungen dasselbe dauerhafte Engineering Object betreffen;
\item ob ein Wert beabsichtigt, realisiert oder beobachtet ist;
\item welches konstruktive Gesetz eine angezeigte Kurve erzeugt hat;
\item welches Wissen anwendbar ist und wer seine Bedeutung verantwortet;
\item ob ein Bewertungsergebnis als Entscheidung angenommen wurde.
\end{itemize}

Dabei handelt es sich nicht um fehlende Attribute mit allgemein
selbstverständlichen Werten. Es sind ingenieurtechnische Unterscheidungen und
Beziehungen, die ein Informationsmodell bewahren muss.

\section{Warum das Alignment die Beziehungen ordnet}

Jedes Informationsmodell wählt ausdrücklich oder implizit einen ordnenden
Bezug. Im Eisenbahnwesen werden viele Fragen entlang oder relativ zu einem
Alignment gestellt: Lage, Richtung, Krümmung, Überhöhung, Geschwindigkeit,
fahrzeugbezogene Größen, Stationierung, Lichträume und die Anordnung
benachbarter Objekte. Das Alignment liefert eine intrinsische longitudinale
Ordnung und einen konstruktiven Krümmungsverlauf, aus denen ausgewählte
Zustände abgeleitet werden können
\cite{Kufver1997MathematicalDescription,Kufver2000RailwayAlignment}.

AIM stellt das Alignment daher in das Zentrum dieser Beziehungen. Dadurch wird
eine Eisenbahn nicht auf ein Objekt reduziert. Brücken, Tunnel, Weichen,
Signale, Stationen, Gelände, Fahrzeuge, Beobachtungen und Entscheidungen
behalten ihre eigenen Rollen. Alignment-zentriert bedeutet, dass Informationen,
deren ingenieurtechnische Bedeutung von der longitudinalen Konstruktion oder
einem gleisrelativen Zustand abhängt, einen ausdrücklichen Bezug besitzen, über
den diese Abhängigkeit interpretiert werden kann.

\begin{definition}[Alignment-Based Information Modelling]
Alignment-Based Information Modelling (AIM) ist ein am Alignment
ausgerichteter Ansatz der Informationsmodellierung, in dem
Alignment-Identität und intrinsische Konstruktion den Bezug für die
Verknüpfung geometrischer, räumlicher, betrieblicher, beobachtender und
bewertender Informationen bilden, ohne das Engineering Object mit einer
einzelnen Repräsentation gleichzusetzen.
\end{definition}

\section{Benachbarte Modellierungsparadigmen}

BIM und AIM müssen nicht miteinander konkurrieren. BIM-Standards stellen
objektorientierte Strukturen und Austauschmechanismen für die gebaute Umwelt
bereit. AIM liefert die Alignment-spezifischen konstruktiven und
zustandsbezogenen Abhängigkeiten, die zur Interpretation von
Alignment-Informationen erforderlich sind. Ein BIM-Arbeitsablauf kann eine aus
AIM abgeleitete Repräsentation tragen; dadurch wird er weder zum Owner der
Alignment-Identität noch des Ingenieurwissens.

Diese Thesis verwendet außerdem \emph{Process Information Modelling} (PIM) als
vorgeschlagenen Oberbegriff für Ansätze, die an Ingenieurprozessen,
Zustandsentwicklung, betrieblichem Verhalten und Lifecycle-Beziehungen statt
ausschließlich an physischen Assets ausgerichtet sind. PIM ist eine Rahmung
dieser Thesis, kein behaupteter externer Standard und keine Voraussetzung für
die AIM-Argumentation.

\section{Die Unterscheidungen, die das Modell bewahren muss}

Das fortlaufende Beispiel verlangt nun sieben Antworten:

\begin{description}[style=nextline]
\item[Objektidentität]
Wodurch betreffen die bearbeiteten und repräsentierten Aufzeichnungen ein
Alignment?
\item[Konstruktive Definition]
Welches geordnete Krümmungsgesetz definiert das Alignment unabhängig von
abgetasteten Koordinaten?
\item[Realisierung]
Wie erhält diese intrinsische Definition metrische Bedeutung, und wie steht
die beabsichtigte Konstruktion zur physischen Infrastruktur?
\item[Repräsentation]
Welche Informationen wurden für diesen GIS-, CAD-, BIM-, numerischen oder
visuellen Zweck ausgewählt?
\item[Beobachtung]
Welche quellengebundene Evidenz berichtet unter welchen Bedingungen und mit
welcher Unsicherheit etwas über das Alignment oder seinen Zustand?
\item[Wissen]
Auf welche Informationen darf innerhalb eines ausdrücklichen Geltungsbereichs
mit benannter Provenienz, Anwendbarkeit und Autorität vertraut werden?
\item[Bewertung und Entscheidung]
Welche Folge wurde berechnet, und welche verantwortete Handlung nimmt eine
Alternative an, lehnt sie ab oder wählt sie aus?
\end{description}

Die Liste wird nicht als auswendig zu lernende Ontologie eingeführt. Jede
Unterscheidung behebt ein bereits an der Krümmungsänderung sichtbares Versagen.
Das nächste Kapitel macht ihre Rollen ausdrücklich, bevor spätere Kapitel die
sicher berechenbaren Beziehungen formalisieren.

\begin{summary}
Repräsentationen machen Ingenieurinformationen für bestimmte Zwecke nutzbar;
der Austausch allein begründet jedoch weder Identität, Zustand,
Anwendbarkeit noch Autorität. AIM ordnet die fehlenden Beziehungen um das
Alignment, sodass eine Ingenieurhandlung über ihre verschiedenen Darstellungen
hinweg nachvollzogen werden kann.
\end{summary}
